\documentclass[aspectratio=169]{beamer}
\usetheme{Madrid}
\usecolortheme{whale}
\usepackage[utf8]{inputenc}
\usepackage[spanish]{babel}
\usepackage{amsmath,amssymb}
\usepackage{graphicx}
\usepackage{booktabs}
\usepackage{listings}
\usepackage{xcolor}

\definecolor{codegreen}{rgb}{0,0.6,0}
\definecolor{codegray}{rgb}{0.5,0.5,0.5}
\definecolor{codepurple}{rgb}{0.58,0,0.82}
\definecolor{backcolour}{rgb}{0.95,0.95,0.92}

\lstdefinestyle{pythonstyle}{
    backgroundcolor=\color{backcolour},
    commentstyle=\color{codegreen},
    keywordstyle=\color{blue},
    numberstyle=\tiny\color{codegray},
    stringstyle=\color{codepurple},
    basicstyle=\ttfamily\footnotesize,
    breakatwhitespace=false,
    breaklines=true,
    captionpos=b,
    keepspaces=true,
    showspaces=false,
    showstringspaces=false,
    showtabs=false,
    tabsize=2
}

\lstset{style=pythonstyle}

\title{Regresión Lineal Múltiple}
\subtitle{Múltiples Predictores}
\author{Métodos de Análisis de Datos 1}
\institute{Universidad Nacional del Sur}
\date{Unidad 3 - Clase 17}

\begin{document}

\frame{\titlepage}

\begin{frame}{Contenidos}
\tableofcontents
\end{frame}

\section{Del modelo simple al múltiple}

\begin{frame}{Motivación}
En la práctica, la respuesta depende de \textbf{múltiples} factores.

\vspace{0.3cm}
\textbf{Ejemplos:}
\begin{itemize}
    \item Precio de casa: metros cuadrados, habitaciones, antigüedad, ubicación
    \item Rendimiento académico: horas de estudio, asistencia, nivel previo
    \item Ventas: publicidad TV, radio, web, estacionalidad
\end{itemize}

\vspace{0.5cm}
\textbf{Ventajas del modelo múltiple:}
\begin{enumerate}
    \item Predicciones más precisas
    \item Controla confusores (variables que afectan tanto $X$ como $Y$)
    \item Permite estudiar efecto de cada variable ajustando por las demás
\end{enumerate}
\end{frame}

\begin{frame}{Modelo de regresión lineal múltiple}
Con $p$ predictores:

\begin{block}{Modelo}
\[
Y_i = \beta_0 + \beta_1 X_{i1} + \beta_2 X_{i2} + \cdots + \beta_p X_{ip} + \epsilon_i
\]
\end{block}

donde:
\begin{itemize}
    \item $Y_i$: variable respuesta
    \item $X_{ij}$: $j$-ésimo predictor para observación $i$
    \item $\beta_0$: intercepto
    \item $\beta_j$: coeficiente del $j$-ésimo predictor
    \item $\epsilon_i \sim N(0, \sigma^2)$: errores independientes
\end{itemize}
\end{frame}

\begin{frame}{Forma matricial}
\begin{block}{Notación matricial}
\[
\mathbf{Y} = \mathbf{X}\boldsymbol{\beta} + \boldsymbol{\epsilon}
\]
\end{block}

\[
\begin{bmatrix} Y_1 \\ Y_2 \\ \vdots \\ Y_n \end{bmatrix} = 
\begin{bmatrix} 
1 & X_{11} & X_{12} & \cdots & X_{1p} \\
1 & X_{21} & X_{22} & \cdots & X_{2p} \\
\vdots & \vdots & \vdots & \ddots & \vdots \\
1 & X_{n1} & X_{n2} & \cdots & X_{np}
\end{bmatrix}
\begin{bmatrix} \beta_0 \\ \beta_1 \\ \vdots \\ \beta_p \end{bmatrix} +
\begin{bmatrix} \epsilon_1 \\ \epsilon_2 \\ \vdots \\ \epsilon_n \end{bmatrix}
\]

\vspace{0.3cm}
\textbf{Dimensiones:} $\mathbf{Y}$ es $n \times 1$, $\mathbf{X}$ es $n \times (p+1)$, $\boldsymbol{\beta}$ es $(p+1) \times 1$
\end{frame}

\begin{frame}{Estimación por MCO}
\begin{block}{Estimador MCO}
\[
\hat{\boldsymbol{\beta}} = (\mathbf{X}^\top \mathbf{X})^{-1} \mathbf{X}^\top \mathbf{Y}
\]
\end{block}

\vspace{0.5cm}
\textbf{Propiedades:} Mismas que en regresión simple
\begin{itemize}
    \item Insesgado
    \item BLUE (mínima varianza entre estimadores lineales insesgados)
    \item Consistente
\end{itemize}

\vspace{0.5cm}
\textbf{Nota:} En la práctica, Python y R hacen estos cálculos automáticamente
\end{frame}

\section{Interpretación}

\begin{frame}{Interpretación de coeficientes}
\textbf{Diferencia clave con regresión simple:}

\begin{block}{Interpretación en modelo múltiple}
$\beta_j$ es el cambio esperado en $Y$ cuando $X_j$ aumenta en 1 unidad, \textbf{manteniendo constantes} las demás variables.
\end{block}

\vspace{0.5cm}
\textbf{Ejemplo:} Precio de casas
\[
\text{Precio} = 50 + 1.2 \times \text{m}^2 + 5 \times \text{habitaciones} - 0.8 \times \text{antigüedad}
\]

\begin{itemize}
    \item $\beta_1 = 1.2$: Por cada m$^2$ adicional, el precio aumenta 1.2 mil \$, \textbf{para igual número de habitaciones y antigüedad}
    \item $\beta_2 = 5$: Cada habitación adicional suma 5 mil \$, \textbf{para igual tamaño y antigüedad}
    \item $\beta_3 = -0.8$: Cada año adicional reduce 0.8 mil \$, \textbf{para igual tamaño y habitaciones}
\end{itemize}
\end{frame}

\begin{frame}{Cuidado con la interpretación}
\begin{alertblock}{Importante}
Los coeficientes solo tienen sentido si los predictores pueden variar independientemente.
\end{alertblock}

\vspace{0.5cm}
\textbf{Problema:} Si hay correlación entre predictores (multicolinealidad), la interpretación se complica.

\vspace{0.3cm}
\textbf{Ejemplo:} Altura y peso están correlacionados
\begin{itemize}
    \item Modelo: IMC $\sim$ altura + peso
    \item ¿Tiene sentido "mantener peso constante y variar altura"?
    \item Los coeficientes individuales pueden ser engañosos
\end{itemize}
\end{frame}

\section{Bondad de ajuste}

\begin{frame}{$R^2$ y $R^2$ ajustado}
\textbf{$R^2$ estándar:}
\[
R^2 = 1 - \frac{\text{SCR}}{\text{SCT}}
\]

\textbf{Problema:} $R^2$ siempre aumenta al agregar variables, aunque sean irrelevantes

\vspace{0.5cm}
\begin{block}{$R^2$ ajustado}
\[
R^2_{\text{adj}} = 1 - \frac{\text{SCR}/(n-p-1)}{\text{SCT}/(n-1)} = 1 - (1-R^2)\frac{n-1}{n-p-1}
\]
\end{block}

\vspace{0.3cm}
\textbf{Ventaja:} Penaliza la complejidad. Puede disminuir si se agregan variables irrelevantes.

\vspace{0.3cm}
\textbf{Uso:} Comparar modelos con distinto número de predictores
\end{frame}

\section{Inferencia}

\begin{frame}{Test F global}
¿Es útil el modelo completo?

\begin{block}{Hipótesis}
\[
H_0: \beta_1 = \beta_2 = \cdots = \beta_p = 0 \quad \text{vs} \quad H_1: \text{al menos un } \beta_j \neq 0
\]
\end{block}

\textbf{Estadístico:}
\[
F = \frac{\text{SCE}/p}{\text{SCR}/(n-p-1)} \sim F_{p, n-p-1}
\]

\vspace{0.3cm}
Si $p$-valor $< \alpha$ $\Rightarrow$ rechazamos $H_0$ (al menos una variable es útil)
\end{frame}

\begin{frame}{Tests t individuales}
Para cada coeficiente:

\begin{block}{Hipótesis}
\[
H_0: \beta_j = 0 \quad \text{vs} \quad H_1: \beta_j \neq 0
\]
\end{block}

\textbf{Estadístico:}
\[
t_j = \frac{\hat{\beta}_j}{\text{SE}(\hat{\beta}_j)} \sim t_{n-p-1}
\]

\vspace{0.5cm}
\textbf{Interpretación:} $X_j$ es significativo \textbf{dado que} las demás variables ya están en el modelo.

\vspace{0.3cm}
\textbf{Cuidado:} Puede ser que $F$ sea significativo pero ningún $t_j$ lo sea (multicolinealidad)
\end{frame}

\section{Multicolinealidad}

\begin{frame}{¿Qué es la multicolinealidad?}
\begin{block}{Definición}
Correlación alta entre dos o más predictores
\end{block}

\vspace{0.3cm}
\textbf{Consecuencias:}
\begin{itemize}
    \item Coeficientes inestables (alta varianza)
    \item Errores estándar grandes
    \item Tests $t$ no significativos, pero test $F$ sí
    \item Interpretación difícil
    \item Predicciones siguen siendo válidas
\end{itemize}

\vspace{0.5cm}
\textbf{No afecta:} $R^2$, predicciones dentro del rango de los datos
\end{frame}

\begin{frame}{Detección: VIF}
\begin{block}{Factor de Inflación de la Varianza}
\[
\text{VIF}_j = \frac{1}{1 - R^2_j}
\]
\end{block}

donde $R^2_j$ es el $R^2$ de la regresión de $X_j$ sobre las demás variables.

\vspace{0.5cm}
\textbf{Interpretación:}
\begin{itemize}
    \item $\text{VIF}_j = 1$: $X_j$ no correlacionado con otras variables
    \item $\text{VIF}_j > 5$: multicolinealidad moderada
    \item $\text{VIF}_j > 10$: multicolinealidad severa
\end{itemize}
\end{frame}

\begin{frame}{Soluciones a la multicolinealidad}
\textbf{Opciones:}

\begin{enumerate}
    \item \textbf{Eliminar variables redundantes}
    \begin{itemize}
        \item Quitar variables muy correlacionadas
    \end{itemize}
    
    \item \textbf{Combinar variables}
    \begin{itemize}
        \item Crear índices o promedios
    \end{itemize}
    
    \item \textbf{Aumentar tamaño muestral}
    \begin{itemize}
        \item Más datos reducen varianza de $\hat{\beta}$
    \end{itemize}
    
    \item \textbf{Regularización}
    \begin{itemize}
        \item Ridge, Lasso (veremos en clases futuras)
    \end{itemize}
    
    \item \textbf{PCA (Análisis de Componentes Principales)}
    \begin{itemize}
        \item Crear predictores no correlacionados
    \end{itemize}
\end{enumerate}
\end{frame}

\section{Implementación}

\begin{frame}[fragile]{Python: Modelo múltiple}
\begin{lstlisting}[language=Python]
from statsmodels.formula.api import ols

# Multiples predictores con formula
modelo = ols('precio ~ metros + habitaciones + antiguedad',
             data=casas).fit()

print(modelo.summary())

# O con matriz de diseño manual
from sklearn.linear_model import LinearRegression
X = casas[['metros', 'habitaciones', 'antiguedad']]
y = casas['precio']
modelo_sk = LinearRegression().fit(X, y)
print(f"Coeficientes: {modelo_sk.coef_}")
print(f"Intercepto: {modelo_sk.intercept_}")
print(f"R^2: {modelo_sk.score(X, y):.4f}")
\end{lstlisting}
\end{frame}

\begin{frame}[fragile]{Python: VIF}
\begin{lstlisting}[language=Python]
from statsmodels.stats.outliers_influence import variance_inflation_factor

# Calcular VIF para cada variable
X = casas[['metros', 'habitaciones', 'antiguedad']]
vif_data = pd.DataFrame()
vif_data["Variable"] = X.columns
vif_data["VIF"] = [variance_inflation_factor(X.values, i) 
                   for i in range(X.shape[1])]
print(vif_data)

#     Variable    VIF
# 0      metros   2.3
# 1 habitaciones  8.5
# 2  antiguedad   1.8
\end{lstlisting}
\end{frame}

\begin{frame}[fragile]{R: Modelo múltiple}
\begin{lstlisting}[language=R]
# Ajustar modelo
modelo <- lm(precio ~ metros + habitaciones + antiguedad,
             data = casas)

summary(modelo)

# VIF
library(car)
vif(modelo)

# Comparar modelos
modelo1 <- lm(precio ~ metros, data = casas)
modelo2 <- lm(precio ~ metros + habitaciones, data = casas)
modelo3 <- lm(precio ~ metros + habitaciones + antiguedad, 
              data = casas)

anova(modelo1, modelo2, modelo3)
\end{lstlisting}
\end{frame}

\section{Resumen}

\begin{frame}{Resumen}
\textbf{Aprendimos:}
\begin{itemize}
    \item Modelo de regresión múltiple: $Y = \beta_0 + \sum \beta_j X_j + \epsilon$
    \item Interpretación: efecto de $X_j$ manteniendo constantes las demás
    \item $R^2$ ajustado para comparar modelos
    \item Test $F$ global y tests $t$ individuales
    \item Multicolinealidad: detección (VIF) y soluciones
    \item Implementación en Python y R
\end{itemize}

\vspace{0.5cm}
\textbf{Próxima clase:} Selección de variables y regresión polinomial
\end{frame}

\end{document}
